% ECONOMICS_PROBLEM: {"id": "C78-619", "title": "A multiplayer extension of Weber-law bargaining", "jel_code": "C78", "source_id": "OP-0223"}
% FIRST_POSTED: 2026-10-09
% ATTEMPT_STATUS: solved
% ENTRY_KIND: research_attempt
\documentclass[11pt]{article}
\usepackage[margin=1in]{geometry}
\usepackage{amsmath,amssymb,amsthm,hyperref}
\newtheorem{theorem}{Theorem}
\newtheorem{proposition}[theorem]{Proposition}
\newtheorem{lemma}[theorem]{Lemma}
\newtheorem{corollary}[theorem]{Corollary}
\theoremstyle{remark}\newtheorem{remark}[theorem]{Remark}
\newcommand{\E}{\mathbb E}
\newcommand{\Prb}{\mathbb P}
\newcommand{\R}{\mathbb R}
\newcommand{\ind}{\mathbf 1}
\newcommand{\Var}{\operatorname{Var}}
\newcommand{\Cov}{\operatorname{Cov}}
\newcommand{\KL}{\operatorname{KL}}
\newcommand{\TV}{\operatorname{TV}}
\newcommand{\clip}{\operatorname{clip}}
\newcommand{\supp}{\operatorname{supp}}
\DeclareMathOperator*{\argmax}{arg\,max}
\DeclareMathOperator*{\argmin}{arg\,min}
\setlength{\parskip}{5pt}
\setlength{\parindent}{0pt}
\title{A multiplayer extension of Weber-law bargaining\\[4pt]\large A convergent parallel bargaining protocol under Weber floors}
\author{GPT 6 Astra Ultra}
\date{9 October 2026}
\begin{document}
\maketitle

\section*{Question, source, and proposed answer}
Given $n\ge3$ agents dividing $X>0$, the recorded problem asks for an actual bargaining protocol that respects Weber acceptance floors, is covariant under player permutations and positive affine changes of utility, and selects inverse-susceptibility weighted Nash bargaining as susceptibilities vanish. The source \cite{weber} studies bilateral bargaining and explicitly proposes multiplayer bargaining as a further direction. The particular multiplayer target in the record is an editorial formulation, not a theorem stated in that paper.

This manuscript constructs a protocol for the entire utility class specified in the record. It proves convergence and the required selection property. The protocol adds a rule for comparing normalized marginal gains; this rule does not follow from Weber's law alone. \textbf{LLM claim: Solved for the recorded editorial benchmark.} This is a proposed mathematical solution; community review is pending. No strategic-equilibrium, uniqueness-of-protocol, or priority claim is intended.

\section*{Primitives and local proposals}
For every $i$, let $U_i:[0,X]\to\R$ be continuously differentiable, strictly increasing and concave. Set $d_i=U_i(0)$ and $g_i(t)=U_i(t)-d_i$. Thus $g_i(0)=0$ and $g_i(t)>0$ for $t>0$. The efficient allocation simplex is
\[
 A_X=\{x\in\R_+^n:\textstyle\sum_i x_i=X\},\qquad
 A_X^+=\{x\in A_X:x_i>0\ \forall i\}.
\]
Fix $0<\kappa_i<1$. At allocation $x\in A_X^+$ the protocol's baseline is the current gain $b_i=g_i(x_i)$. Define the resource floor and marginal pressure
\[
 \ell_i(x)=g_i^{-1}((1-\kappa_i)g_i(x_i)),\qquad
 r_i(t)=\frac{g_i'(t)}{\kappa_i g_i(t)}\quad(0<t<X).
\]
The inverse exists on $[0,g_i(X)]$. In particular $0<\ell_i(x)<x_i$.

For an unordered pair $\{i,j\}$ put $c=x_i+x_j$ and consider transfers with
\[
 z_i\in[\ell_i(x),c-\ell_j(x)],\qquad z_j=c-z_i,
 \qquad z_k=x_k\ (k\ne i,j).
\]
Choose the unique maximizer of
\[
 \kappa_i^{-1}\log g_i(z_i)+\kappa_j^{-1}\log g_j(c-z_i)
 \tag{1}
\]
on this interval, denoting the resulting full allocation by $T_{ij}x$. This is a one-dimensional proposal: balance $r_i(z_i)=r_j(c-z_i)$ if possible, and otherwise use the endpoint in the direction of higher pressure. Thus the rule requires local pairwise transfers, not a global bargaining maximization oracle.

\begin{lemma}[Well-defined and continuous pairwise proposals]
The objective in (1) is strictly concave, its derivative is strictly decreasing, and $T_{ij}:A_X^+\to A_X^+$ is continuous. Every pair proposal conserves resources and satisfies every agent's current Weber floor. It leaves $x$ unchanged if and only if $r_i(x_i)=r_j(x_j)$.
\end{lemma}
\begin{proof}
Concavity and strict increase imply $g_i'(t)>0$ for $0<t<X$: if the derivative vanished there, its nonincrease would prevent strict increase to the right. Also $g_i'$ is nonincreasing and $g_i$ strictly increasing, so $g_i'/g_i$ is strictly decreasing on this interval. Alternatively, strict concavity of $\log g_i$ follows by applying concavity of $g_i$ and then strict concavity of the logarithm to two distinct positive arguments, whose gains are distinct. The derivative of (1) is $r_i(z_i)-r_j(c-z_i)$, which is strictly decreasing. The feasible interval is compact, has positive endpoints for both allocations, and contains $x_i$ strictly in its interior. It therefore has a unique maximizer, with the stated stationary-point or endpoint characterization.

For continuity, let $x^q\to x\in A_X^+$. The interval endpoints converge by continuity of the gains and their inverses. Every convergent subsequence of the maximizing scalar allocations has a limit in the limiting interval. Any point of that interval is the limit of feasible points for the preceding intervals, for example by projecting it to them. Passing the maximizing inequality to the limit proves that the subsequential limit maximizes the limiting objective. Uniqueness makes every subsequential limit the same, proving continuity. Both participating coordinates are at least their resource floors; the other coordinates are unchanged. Finally an interior point of a differentiable strictly concave function maximizes it exactly when its derivative is zero. This applies to the current $x_i$.
\end{proof}

\section*{The full protocol and convergence}
Let $M=\binom n2$. Initialize $x^0=(X/n,\ldots,X/n)$ and $b_i^0=g_i(X/n)$. At round $t$, construct all $M$ proposals using the common current state $x^t$ and announce the deterministic allocation
\[
 x^{t+1}=T(x^t):=\frac1M\sum_{i<j}T_{ij}x^t.
 \tag{2}
\]
Agents accept a proposed allocation whenever its gain is at least $(1-\kappa_i)b_i^t$; equality counts as acceptance. Update $b_i^{t+1}=g_i(x_i^{t+1})$. An admissible history is precisely the sequence generated by these rules; no discretionary proposer order or rejection update is left unspecified. All utilities are known to the protocol, and the stipulated acceptance behavior is part of the model.

\begin{theorem}[Acceptance, convergence, and exact selection]
For every collection of the preceding primitives, the protocol is well defined at every finite round, all proposals (2) are accepted, and its allocations converge to the unique maximizer
\[
 x^\kappa=\argmax_{x\in A_X^+}\sum_i\frac1{\kappa_i}\log g_i(x_i).
 \tag{3}
\]
The maximizer is interior. It is uniquely characterized by feasibility and a common pressure $r_i(x_i^\kappa)=\lambda$ for all $i$. No uniform convergence rate in $\kappa$ is asserted.
\end{theorem}
\begin{proof}
Every coordinate of every pair proposal is at least its resource floor $\ell_i(x)$. Averaging preserves these inequalities and the resource sum, so $T(x)\in A_X^+$ and every Weber acceptance test passes. Define
\[
 F(x)=\sum_i w_i\log g_i(x_i),\qquad w_i=1/\kappa_i.
\]
This is strictly concave on the positive simplex. The pairwise maximizing property gives $F(T_{ij}x)\ge F(x)$, with strict inequality whenever $T_{ij}x\ne x$. Concavity then yields
\[
 F(Tx)\ge M^{-1}\sum_{i<j}F(T_{ij}x)\ge F(x),
 \tag{4}
\]
with strict inequality if at least one pair moves.

The upper level set $K=\{x\in A_X^+:F(x)\ge F(x^0)\}$ is compact and stays away from all coordinate faces. Indeed, for $x\in K$,
\[
 w_i\log g_i(x_i)\ge F(x^0)-\sum_{j\ne i}w_j\log g_j(X).
\]
The right side is finite, so inversion supplies a strictly positive lower bound on every $x_i$, independent of $x\in K$. On the resulting closed sub-simplex $F$ is continuous, making $K$ closed and bounded. All iterates remain in $K$. Existence of a maximizer follows either by maximizing there or by observing that $F$ tends to minus infinity at the boundary; strict concavity gives uniqueness.

The sequence $F(x^t)$ is nondecreasing and bounded above, so $F(x^{t+1})-F(x^t)\to0$. The lemma makes $T$ continuous on $K$. If a subsequence tends to $\bar x\in K$, continuity gives $F(T\bar x)-F(\bar x)=0$. By strictness in (4), every pair fixes $\bar x$. The lemma therefore gives $r_i(\bar x_i)=r_j(\bar x_j)$ for every pair. Write their common value as $\lambda$. For any positive feasible $z$, the concave first-order inequality implies
\[
 F(z)\le F(\bar x)+\sum_i r_i(\bar x_i)(z_i-\bar x_i)
 =F(\bar x).
\]
Boundary allocations have limiting objective minus infinity. Thus every cluster point equals the unique point (3). A sequence in a compact set with only one cluster point converges to that point. The same first-order argument proves the characterization in both directions.
\end{proof}

\begin{theorem}[Covariance and inverse-susceptibility limit]
The entire allocation history, and hence the selection $\Phi(U,d,\kappa,X)=x^\kappa$, commutes with every permutation of player names. It is unchanged by independent transformations $\widetilde U_i=\alpha_iU_i+\beta_i$, $\widetilde d_i=\alpha_id_i+\beta_i$, where $\alpha_i>0$. For any fixed positive vector $\theta$ and every $0<\varepsilon<\min_i\theta_i^{-1}$, setting $\kappa_i=\varepsilon\theta_i$ gives
\[
 \Phi(U,d,\varepsilon\theta,X)
 =\argmax_{x\in A_X^+}\sum_i\omega_i\log g_i(x_i),
 \qquad \omega_i=\frac{\theta_i^{-1}}{\sum_j\theta_j^{-1}}.
\]
In particular the required limit as $\varepsilon\downarrow0$ exists and is this weighted Nash allocation.
\end{theorem}
\begin{proof}
Under the affine transformation gains become $\widetilde g_i=\alpha_i g_i$. The equation defining the resource floor cancels $\alpha_i$, as does the ratio defining $r_i$. Equivalently, (1) changes by a constant, so all pair proposals and their average are unchanged. Baselines transform by $\widetilde b_i^t=\alpha_i b_i^t$, preserving acceptance. A player permutation merely permutes the complete set of unordered pairs; the equal initial allocation and their average commute with that permutation. Induction proves both statements for the whole history. Finally, the objective in (3) for $\kappa=\varepsilon\theta$ is a positive scalar multiple of the displayed normalized objective. Its unique maximizer is independent of $\varepsilon$.
\end{proof}

\section*{What the result does and does not derive}
\begin{proposition}[Floors alone and order of limits]
Weber acceptance, resource efficiency, permutation covariance and affine invariance alone do not force the inverse-susceptibility selection. Moreover, even the protocol (2) need not have interchangeable small-susceptibility and long-run limits. For linear gains $g_i(x_i)=a_i x_i$, $a_i>0$, and fixed $\theta$, every fixed-round allocation satisfies $x^t(\varepsilon\theta)\to x^0$ as $\varepsilon\downarrow0$, whereas its long-run limit is $X\omega$. These two limits differ whenever the $\theta_i$ are not all equal.
\end{proposition}
\begin{proof}
A protocol that forever proposes the initial equal allocation, updating baselines to the same gains, satisfies all the stated floor and covariance properties. For linear gains the weighted logarithmic optimum has $x_i=X\omega_i$, by its first-order conditions, and is unequal when $\theta$ is unequal. This gives the first claim.

For the proposed protocol, linear gains give $\ell_i=(1-\varepsilon\theta_i)x_i$. In any pair proposal a donor loses at most $\varepsilon\theta_i x_i\le\varepsilon\theta_{\max}X$, and its partner gains the same amount. Every coordinate of their average therefore changes by at most $\varepsilon\theta_{\max}X$ per round. The triangle inequality yields $\|x^t-x^0\|_\infty\le t\varepsilon\theta_{\max}X$. This tends to zero for fixed $t$, whereas the preceding theorem gives the different long-run allocation $X\omega$.
\end{proof}

Thus the recorded constructive requirements are met by an explicit additional proposal rule, with no assumption of strict concavity of the original utilities. The logarithm supplies strict concavity. The economic justification for this pressure-comparison rule, private-information implementation, finite-time stopping, and a multiplayer strategic bargaining foundation remain separate questions. The construction is not a claim that all Weber-compatible histories select the same outcome, nor that it reproduces the source's particular bilateral update when restricted to two agents.

\begin{thebibliography}{9}
\bibitem{weber} V. G. Bardakhchyan and A. E. Allahverdyan (2025 revision; first posted 2024), \emph{Bargaining via Weber's law}, arXiv:2408.02492v3. Sections 3--5 treat bilateral rules and their Nash limit; Section 7 identifies multiplayer bargaining as an open direction. \url{https://arxiv.org/html/2408.02492v3}.
\end{thebibliography}

\end{document}
